\color{unidarkgreen}{\subsection{Комбинированный пусковой орган напряжения}}
\color{black}

\begin{enumerate}

\item
Комбинированный пусковой орган по напряжению своей стороны (КПОН СШ) и низкой стороны (КПОН НН) используется для повышения чувствительности ступеней МТЗ.  

\item
Уставки КПОН представлены в таблице \ref{app_set:kpon}. На рисунке \ref{pic_kpon:KPON} приведен алгоритм работы КПОН. 

\begin{figure}[H]
\centering
\includegraphics[width=0.8\textwidth,height=0.8\textheight,keepaspectratio]{\fbpath/07.02 РЗАТ/051.0371 КПОН/img/KPON.pdf}
\captionof{figure}{Функциональная схема логики КПОН}\label{pic_kpon:KPON}
\end{figure}

\item
Цепи напряжения подключаются к ТН, установленным на шинах низшего напряжения.

\item
Внутренний комбинированный пусковой орган напряжения низкой стороны (КПОН НН), контролирует минимальное линейное напряжение и максимальное напряжение обратной последовательности от ТН, установленного на ошиновке НН или на вводе секции. Порог срабатывания пуска по минимальному линейному напряжению задается уставкой \textbf{«Uл КПОН НН»}, порог срабатывания пуска по максимальному напряжению обратной последовательности задается уставкой \textbf{«U2 КПОН НН»}. Ввод в работу КПОН НН осуществляется с помощью программного переключателя \textbf{«КПОН НН»}. 

\item
Выбор пуска по напряжению КПОН НН осуществляется с помощью программного переключателя \textbf{«Тип КПОН НН»}: 
\begin{itemize}
\item \textit{«Umin, U2»} -- Пуск по минимальному линейному напряжению или максимальному напряжению обратной последовательности;
\item \textit{«Umin»} -- Пуск по минимальному линейному напряжению;
\item \textit{«Внеш»} -- Пуск по внешнему входному сигналу \textit{«Внеш. КПОН»} соответствующей стороны.
\end{itemize}

\item
Предусмотрена возможность контроля неисправности цепей напряжения стороны НН по внутренней логике КЦН НН. При появлении неисправности цепей напряжения, действие логики КПОН НН задается программным переключателем \textbf{«Неиспр ТН НН»}:
\begin{itemize}
\item \textit{«Выведено»} -- Неисправность в цепях ТН не вызывает изменения логики КПОН;
\item \textit{«Пуск КПОН»} -- При появлении неисправности в цепях ТН, КПОН всегда разрешен;
\item \textit{«Блок КПОН»} -- При появлении неисправности в цепях ТН, КПОН блокируется до исчезновения неисправности ТН.
\end{itemize}

\item
Оперативный ввод/вывод КПОН НН производится с помощью БУ (блоков управления функциями) \textit{«Вывод КПОН НН»}. 

\item
Внутренний комбинированный пусковой орган напряжения своей стороны (КПОН СШ), контролирует минимальное линейное напряжение и максимальное напряжение обратной последовательности от ТН, установленного на системе шин своей стороны. Порог срабатывания пуска по минимальному линейному напряжению задается уставкой \textbf{«Uл КПОН СН»}, порог срабатывания пуска по максимальному напряжению обратной последовательности задается уставкой \textbf{«U2 КПОН СШ»}. Ввод в работу КПОН СШ осуществляется с помощью программного переключателя \textbf{«КПОН СШ»}.

\item
Выбор пуска по напряжению КПОН СШ осуществляется с помощью программного перключателя \textbf{«Тип КПОН СШ»}:
\begin{itemize}
\item \textit{«Umin, U2»} -- Пуск по минимальному линейному напряжению или максимальному напряжению обратной последовательности;
\item \textit{«Umin»} -- Пуск по минимальному линейному напряжению.
\end{itemize}

\item
Предусмотрена возможность контроля неисправности цепей напряжения своей стороны по внутренней логике БНН. При появлении неисправности в цепях напряжения, действие логики КПОН СШ задается программным переключателем \textbf{«Неиспр ТН СШ»}:
\begin{itemize}
\item \textit{«Выведено»} -- Неисправность в цепях ТН не вызывает изменения логики КПОН;
\item \textit{«Пуск КПОН»} -- При появлении неисправности в цепях ТН, КПОН всегда разрешен;
\item \textit{«Блок КПОН»} -- При появлении неисправности в цепях ТН, КПОН блокируется до исчезновения неисправности ТН.
\end{itemize}

\item
Оперативный ввод/вывод КПОН СШ производится с помощью БУ (блоков управления функциями) \textit{«Вывод КПОН СШ»}. 

\end{enumerate}